\documentclass[11pt,a4paper]{article}

\usepackage[T1]{fontenc}
\usepackage[utf8]{inputenc}
\usepackage{lmodern}
\usepackage{textcomp}
\usepackage[margin=2.2cm]{geometry}
\usepackage{xcolor}
\usepackage{array}
\usepackage{tabularx}
\usepackage{longtable}
\usepackage{booktabs}
\usepackage{enumitem}
\usepackage{fancyhdr}
\usepackage{titlesec}
\usepackage{listings}
\usepackage[colorlinks=true,linkcolor=maroon,urlcolor=maroon,citecolor=maroon]{hyperref}

\definecolor{maroon}{HTML}{8C1C2B}
\definecolor{codebg}{HTML}{F5F5F5}

\newcommand{\degC}{\textdegree{}C}
\newcommand{\us}{\textmu s}
\newcommand{\ttu}{\textunderscore}

\titleformat{\section}{\Large\bfseries\color{maroon}}{\thesection}{0.6em}{}[\vspace{-0.4em}\color{maroon}\titlerule]
\titleformat{\subsection}{\large\bfseries}{\thesubsection}{0.6em}{}
\titlespacing*{\section}{0pt}{1.6em}{0.8em}

\pagestyle{fancy}
\fancyhf{}
\lhead{\small Vehicle Event Logging and Response System}
\rhead{\small 23CSE304 -- Embedded System}
\cfoot{\small\thepage}
\renewcommand{\headrulewidth}{0.4pt}

\setlength{\parindent}{0pt}
\setlength{\parskip}{0.5em}
\renewcommand{\arraystretch}{1.25}
\setlist{itemsep=0.1em,topsep=0.3em}

\lstset{
  basicstyle=\ttfamily\footnotesize,
  backgroundcolor=\color{codebg},
  frame=single,
  rulecolor=\color{black!25},
  breaklines=true,
  columns=fullflexible,
  keepspaces=true,
  xleftmargin=0.3em,
  xrightmargin=0.3em,
  aboveskip=0.6em,
  belowskip=0.6em
}

\newcolumntype{L}[1]{>{\raggedright\arraybackslash}p{#1}}

\begin{document}

%% ---------------------------------------------------------------- Title
\begin{titlepage}
\centering
\vspace*{1.5cm}
{\Huge\bfseries\color{maroon} Vehicle Event Logging\\[0.2em] and Response System\par}
\vspace{0.8cm}
{\Large Case Study Report\par}
\vspace{0.3cm}
{\large STM32F411CEU6 (ARM Cortex-M4) -- Bare-Metal Embedded C\par}
\vspace{1.6cm}
\rule{0.8\textwidth}{0.6pt}\\[1em]
{\Large\bfseries Academic Details and Acknowledgement\par}
\vspace{1em}
\begin{flushleft}
\textbf{Student Name(s) and Registration Number(s):}
\begin{center}
\begin{tabular}{ll}
\toprule
\textbf{Registration No.} & \textbf{Name}\\
\midrule
CB.SC.U4CSE24502 & Aniruth Karthik\\
CB.SC.U4CSE24513 & Jaikrrish S\\
CB.SC.U4CSE24521 & Krishanth K\\
CB.SC.U4CSE24550 & Srisanth P A\\
\bottomrule
\end{tabular}
\end{center}
\textbf{Mentor:} Dr.\ A.\ Baskar\\[0.3em]
\textbf{Course:} 23CSE304 -- Embedded System\\[0.3em]
\textbf{Department:} Department of Computer Science and Engineering\\[0.3em]
\textbf{Institution:} School of Computing, Amrita Vishwa Vidyapeetham, Coimbatore, India.
\end{flushleft}
\vfill
{\large Team 11}
\end{titlepage}

\tableofcontents
\newpage

%% ---------------------------------------------------------------- 1
\section{Abstract}
This project implements a \textbf{Vehicle Event Logging and Response System} on an STM32F411CEU6 ``Black Pill'' microcontroller. Three sensors are monitored continuously: an \textbf{MPU6050} accelerometer/gyroscope (impact and motion, over I\textsuperscript{2}C), an \textbf{LM35} analog temperature sensor (over the ADC), and an \textbf{HC-SR04} ultrasonic ranger (obstacle distance, using GPIO and cycle-accurate timing). Readings are classified in firmware as \emph{Normal}, \emph{Warning}, \emph{Critical} or \emph{Invalid} against configurable thresholds. Critical events are transmitted as structured ASCII alerts over UART at 115200 baud and are also indicated locally by an LED. The firmware is written at the \textbf{register level} using only CMSIS device headers; it uses no HAL and no RTOS.

%% ---------------------------------------------------------------- 2
\section{Introduction and Problem Statement}
A vehicle can suffer sudden impact, overheating, or close-range obstacles with no automated mechanism to detect, record or report the event. Without continuous monitoring, such events may go unnoticed or be reported too late.

This system addresses that gap by:
\begin{itemize}
  \item continuously sampling vehicle-mounted sensors,
  \item detecting abnormal conditions by comparing readings against defined thresholds,
  \item giving an immediate local indication (LED), and
  \item sending a structured alert message to a PC terminal, which stands in for notification of emergency services.
\end{itemize}

\textit{Demonstration note:} Alerts are received on a PC serial terminal (Docklight / PuTTY / CoolTerm) through a USB-to-UART adapter. \texttt{[ALERT]} messages represent the emergency-notification channel.

%% ---------------------------------------------------------------- 3
\section{Objectives and Scope}
\subsection{Objectives}
\begin{tabularx}{\textwidth}{@{}l L{3.6cm} X@{}}
\toprule
\textbf{\#} & \textbf{Objective} & \textbf{Realisation}\\
\midrule
O1 & Temperature monitoring & LM35 + ADC1, 32-sample averaging, over-temperature alert\\
O2 & Impact / motion detection & MPU6050 over I\textsuperscript{2}C, baseline subtraction, threshold on peak axis deviation\\
O3 & Obstacle detection & HC-SR04 trigger/echo timing, distance in cm, proximity alert\\
O4 & Local response & PC13 LED: 1~Hz heartbeat when normal, 10~Hz strobe when critical\\
O5 & Remote notification & Structured ASCII alerts and telemetry over UART at 115200 8-N-1\\
O6 & Asynchronous emergency input & EXTI interrupt on PB12 (manual crash/SOS trigger)\\
O7 & Register-level implementation & No HAL; direct access to RCC, GPIO, ADC, USART, EXTI, SYSCFG and DWT registers\\
\bottomrule
\end{tabularx}

\subsection{In Scope (as implemented)}
\begin{itemize}
  \item MPU6050, LM35 and HC-SR04 sensing
  \item Threshold-based classification in firmware
  \item UART alert and telemetry transmission
  \item LED indication and the PB12 hardware emergency interrupt
\end{itemize}

\subsection{Features Not Implemented}
The original project plan describes a wider system than the one built. Features that were planned but not implemented are listed in Section~\ref{sec:outofscope}.

%% ---------------------------------------------------------------- 4
\section{System Overview and Architecture}
\subsection{Block Diagram}
\begin{lstlisting}
              +-------------------------------------------------------+
              |            STM32F411CEU6 (Cortex-M4 @ 16 MHz HSI)     |
 LM35  -----> | PA0 / ADC1_IN0 --+                                    |
              |                  |                                    |
 HC-SR04 ---> | PB0 TRIG (out)   +--> Sensor drivers --> Classifier   |
              | PB1 ECHO (in)    |    (temp / dist / motion)          |
              |                  |             |                      |
 MPU6050 ---> | PB6 SCL / PB7 SDA (bit-banged I2C)                    |
              |                                v                      |
 SOS button-> | PB12 / EXTI12 --ISR--> Event Manager                  |
              |                      (alerts, telemetry, LED)         |
              +-------------------------|------------|----------------+
                                        v            v
                       UART TX (PA9/PA2/PB10/PB4)  PC13 LED
                                        |
                                USB-TTL --> PC terminal (Docklight)
\end{lstlisting}

\subsection{Processing Pipeline}
\begin{tabularx}{\textwidth}{@{}l L{2.6cm} X@{}}
\toprule
\textbf{Step} & \textbf{Stage} & \textbf{Mechanism}\\
\midrule
1 & Sense & ADC conversion (LM35); echo-pulse timing (HC-SR04)\\
2 & Acquire & 14-byte burst read of the MPU6050 over I\textsuperscript{2}C\\
3 & Process & Averaging, scaling, baseline subtraction\\
4 & Detect & Threshold comparison $\rightarrow$ \texttt{EventState}\\
5 & Interrupt & PB12 falling-edge EXTI sets a critical motion state immediately\\
6 & Respond & LED heartbeat or strobe on PC13\\
7 & Communicate & ASCII alerts and \texttt{[DATA]} telemetry over UART\\
\bottomrule
\end{tabularx}

\subsection{Main-Loop Timing}
The application is a \textbf{cooperative super-loop} with a nominal 100~ms period (PRD requirement NF-01 allows up to 200~ms).
\begin{lstlisting}
INIT  clock -> DWT -> GPIO -> EXTI -> UART -> ADC -> banner
      -> MPU6050 init + baseline
LOOP  (every ~100 ms)
        1. lm35_read_and_classify()
        2. hcsr04_read_and_classify()
        3. mpu6050_read_and_classify()  (or re-init if offline)
        4. transmit_alerts_and_telemetry()
        5. delay_ms(100)
\end{lstlisting}
Critical alerts are rate-limited to every 5th loop ($\approx$0.5~s) and the normal heartbeat to every 10th loop ($\approx$1~s), so a persistent fault does not flood the UART.

%% ---------------------------------------------------------------- 5
\section{Hardware Design}
\subsection{Components}
\begin{tabularx}{\textwidth}{@{}L{4.2cm} X L{2.2cm}@{}}
\toprule
\textbf{Component} & \textbf{Role} & \textbf{Interface}\\
\midrule
STM32F411CEU6 Black Pill & Main controller; Cortex-M4F, run at 16~MHz HSI & ---\\
MPU6050 (GY-521) & 3-axis accelerometer, 3-axis gyroscope, die temperature & I\textsuperscript{2}C\\
LM35 & Analog temperature sensor, 10~mV/\degC & ADC\\
HC-SR04 & Ultrasonic distance sensor, 2--400~cm & GPIO + timing\\
USB-to-UART adapter (CP2102/CH340) & Bridge to the PC terminal & UART\\
On-board LED (PC13, active LOW) & Local status indication & GPIO\\
Push-button (PB12 to GND) & Manual emergency / SOS input & GPIO + EXTI\\
\bottomrule
\end{tabularx}

\subsection{Pin Map (from \texttt{config.h})}
\begin{tabularx}{\textwidth}{@{}L{3.2cm} L{1.8cm} L{3.8cm} X@{}}
\toprule
\textbf{Signal} & \textbf{Pin} & \textbf{Function} & \textbf{Notes}\\
\midrule
LM35 output & PA0 & ADC1\_IN0 & Analog mode, no pull\\
HC-SR04 TRIG & PB0 & GPIO output & Push-pull, idle LOW\\
HC-SR04 ECHO & PB1 & GPIO input, pull-down & 5~V-tolerant pin, wired directly\\
MPU6050 SCL & PB6 & Open-drain, pull-up & Bit-banged I\textsuperscript{2}C\\
MPU6050 SDA & PB7 & Open-drain, pull-up & Bit-banged I\textsuperscript{2}C; AD0 to GND gives 0x68\\
UART TX (primary) & PA9 & USART1 TX, AF7 & Hardware UART\\
UART TX (mirror) & PA2 & USART2 TX, AF7 & Hardware UART\\
UART TX (mirror) & PB10, PB4 & GPIO, bit-banged & 139 cycles per bit\\
Status LED & PC13 & GPIO output & Active LOW\\
Emergency input & PB12 & EXTI12, falling edge, pull-up & NVIC priority 2\\
\bottomrule
\end{tabularx}

All pin numbers and thresholds are centralised in \texttt{config.h}; other code refers to them only through macros (PRD Modularity Rule~1).

\subsection{Clocking}
The firmware runs from the internal \textbf{16~MHz HSI} oscillator with all bus prescalers at 1 and zero flash wait states. The \textbf{DWT cycle counter} (\texttt{CYCCNT}) is the free-running time base. It gives cycle-exact delays (\texttt{delay\ttu us}, \texttt{delay\ttu ms}) and pulse-width measurement without a hardware timer.

%% ---------------------------------------------------------------- 6
\section{Embedded Concepts Applied}
\begin{tabularx}{\textwidth}{@{}L{3.4cm} L{3.6cm} X@{}}
\toprule
\textbf{Concept} & \textbf{Where used} & \textbf{What it demonstrates}\\
\midrule
Cortex-M4 core peripherals & DWT, CoreDebug, NVIC & Debug-unit cycle counter as a time base; interrupt priority control\\
Memory-mapped register I/O & All drivers & Direct MODER, AFR, BSRR, CR1/CR2, SR/DR access\\
GPIO & Trigger, echo, LED, button, bit-banged buses & Input, output, alternate and analog modes; atomic set/reset via BSRR\\
ADC & LM35 & 12-bit single conversion, 480-cycle sample time, software oversampling\\
Serial buses & MPU6050 (I\textsuperscript{2}C), alert link (UART) & Master protocol implementation, ACK handling, bus recovery\\
Interrupts (EXTI/NVIC) & PB12 emergency input & Short ISR, flag passing to the main loop, write-buffer drain\\
Timing & DWT delay and echo capture & Deterministic, hardware-independent timing\\
Fault tolerance & MPU6050 re-init, HC-SR04 timeouts & Graceful degradation; invalid readings never raise alerts\\
Observability & \texttt{volatile} globals & Live inspection in the Keil Watch window\\
\bottomrule
\end{tabularx}

%% ---------------------------------------------------------------- 7
\section{Firmware Design}
The main firmware is a single translation unit, \texttt{main.c}, with configuration in \texttt{config.h}.

\subsection{Clock and Timing}
\texttt{clock\ttu init()} enables HSI, selects it as SYSCLK and clears the prescalers. \texttt{dwt\ttu init()} enables the trace block and starts \texttt{CYCCNT}. \texttt{delay\ttu us()} busy-waits on \texttt{CYCCNT} using unsigned subtraction, so it stays correct across counter wrap-around.

\subsection{LM35 Temperature}
\begin{itemize}
  \item ADC1 is configured for 12-bit single conversion on channel~0 with the longest sample time (480 cycles), because the LM35 has a high source impedance.
  \item \textbf{32 conversions} are averaged per reading to suppress noise.
  \item Conversion: $\mathrm{mV} = \mathrm{raw}\times 3300/4095$, then $T = \mathrm{mV}/10$ (LM35 gives 10~mV/\degC).
  \item Classification uses the rounded integer \degC{} value.
\end{itemize}

\subsection{HC-SR04 Distance}
\begin{enumerate}
  \item Drive TRIG HIGH for 15~\us{} (the sensor needs at least 10~\us), then LOW.
  \item Wait for ECHO to rise, with a \textbf{35~ms timeout}. No echo means no target or no sensor.
  \item Time how long ECHO stays HIGH using \texttt{CYCCNT}, with a \textbf{40~ms} ceiling.
  \item $d_{\mathrm{cm}} = t_{\mathrm{us}}/58$ (sound travels about 58~\us{} per cm for the round trip).
  \item Only values in \textbf{2--400~cm} are accepted. Anything else, and any timeout, gives \texttt{STATE\ttu INVALID} with no alert (PRD NF-07).
\end{enumerate}

\subsection{Bit-Banged I\textsuperscript{2}C}
I\textsuperscript{2}C is implemented in software on PB6/PB7 at roughly \textbf{50~kHz} (10~\us{} half-period). It provides:
\begin{itemize}
  \item START/STOP generation and byte write/read with ACK/NACK,
  \item \textbf{clock-stretching support}: \texttt{i2c\ttu scl\ttu hi()} waits for the line to actually go high, with a timeout,
  \item \textbf{bus clear}: \texttt{i2c\ttu bus\ttu clear()} sends nine clock pulses plus a STOP to release a slave holding SDA low,
  \item register write, single-register read and burst read.
\end{itemize}

\subsection{MPU6050 Driver}
\begin{itemize}
  \item \textbf{Auto-probe} of addresses 0x68 and 0x69, validated by reading \texttt{WHO\ttu AM\ttu I} (0x75) and rejecting 0x00 and 0xFF (a floating bus).
  \item \textbf{Initialisation:} wake from sleep, PLL clock with X-gyro reference, enable all axes, \texttt{SMPLRT\ttu DIV}=0x07 (125~Hz), \texttt{CONFIG}=0x06 (DLPF $\approx$5~Hz), accelerometer $\pm$2~g, gyroscope $\pm$250~\textdegree/s, interrupts disabled.
  \item \textbf{Burst read:} all 14 bytes (accelerometer, temperature, gyroscope) in one transaction from 0x3B, combined into signed 16-bit values.
  \item \textbf{Baseline calibration:} 32 stationary samples, 10~ms apart, are averaged per axis. Live samples are compared against this baseline, so gravity and mounting orientation do not trigger false events.
  \item \textbf{Scaling:} $\pm$2~g gives 16384~LSB/g and $\pm$250~\textdegree/s gives 131~LSB/(\textdegree/s). The impact metric is $\max(|\Delta a_x|,|\Delta a_y|,|\Delta a_z|)\times 100/16384$, i.e.\ g$\times$100 in integer arithmetic.
  \item \textbf{Self-healing:} on a read failure the state becomes \texttt{STATE\ttu INVALID}, the bus is cleared and the sensor is re-initialised. If startup detection failed, the main loop retries \texttt{mpu6050\ttu init()} every cycle.
\end{itemize}

\subsection{Emergency Interrupt}
PB12 is an input with pull-up, routed through \texttt{SYSCFG->EXTICR[3]} to EXTI12 and triggered on the \textbf{falling edge}. The ISR is short and:
\begin{enumerate}
  \item clears the pending bit (write-1-to-clear) and reads \texttt{EXTI->PR} back to drain the write buffer, preventing spurious re-entry,
  \item sets \texttt{emergency\ttu exti\ttu triggered},
  \item forces \texttt{state\ttu motion = STATE\ttu CRITICAL}, increments the alert counter and toggles the LED.
\end{enumerate}
The ISR does no I\textsuperscript{2}C and no UART (PRD NF-05); the main loop prints the alert.

\subsection{UART Driver}
Each character is sent on \textbf{four} outputs at 115200 baud, 8-N-1:

\begin{tabularx}{\textwidth}{@{}L{2.4cm} X@{}}
\toprule
\textbf{Output} & \textbf{Method}\\
\midrule
PA9 & USART1 (hardware), BRR = 139 (16~MHz / 115200)\\
PA2 & USART2 (hardware), BRR = 139\\
PB10, PB4 & Software UART: start bit, 8 data bits LSB-first, stop bit, each 139 cycles long, timed from \texttt{CYCCNT}\\
\bottomrule
\end{tabularx}

The mirrored outputs let the team connect whichever USB-TTL adapter or header pin is available without reflashing. Hardware TXE waits have a 5000-iteration timeout so a disconnected port cannot hang the loop. \texttt{uart\ttu send\ttu string()} expands \texttt{\textbackslash n} to \texttt{\textbackslash r\textbackslash n}.

\subsection{Event Manager}
\begin{itemize}
  \item Emits \texttt{[ALERT]} lines for each sensor in \texttt{STATE\ttu CRITICAL} (every 5th loop).
  \item Emits the \texttt{[INFO] SYSTEM | NORMAL} heartbeat plus a \texttt{[DATA]} telemetry line when nothing is critical (every 10th loop).
  \item Drives the LED: 10~Hz toggle on any critical state, 1~Hz heartbeat otherwise.
\end{itemize}

\subsection{Debug Observability}
Sensor values, raw counts, baselines, states and counters (\texttt{temperature}, \texttt{adc\ttu mv}, \texttt{distance}, \texttt{echo\ttu pulse\ttu us}, \texttt{accel\ttu x/y/z}, \texttt{max\ttu delta\ttu g\ttu x100}, \texttt{critical\ttu alert\ttu count}, \texttt{mpu\ttu id}, \dots) are \texttt{volatile} globals that can be watched live in Keil \textmu Vision without UART.

%% ---------------------------------------------------------------- 8
\section{Event Classification Logic}
All thresholds live in \texttt{config.h}:

\begin{tabularx}{\textwidth}{@{}L{3.4cm} L{2.8cm} L{2.8cm} X@{}}
\toprule
\textbf{Sensor} & \textbf{Warning} & \textbf{Critical} & \textbf{Invalid}\\
\midrule
Temperature (LM35) & $T\ge 45$~\degC & $T\ge 60$~\degC & ---\\
Distance (HC-SR04) & $d\le 35$~cm & $d\le 20$~cm & Timeout, or $d$ outside 2--400~cm\\
Motion (MPU6050) & $\max\Delta\ge 0.80$~g & $\max\Delta\ge 1.50$~g & I\textsuperscript{2}C read failure\\
\bottomrule
\end{tabularx}

\begin{lstlisting}[language=C]
typedef enum {
    STATE_NORMAL   = 0,   /* within safe range             */
    STATE_WARNING  = 1,   /* approaching threshold         */
    STATE_CRITICAL = 2,   /* threshold exceeded -> alert   */
    STATE_INVALID  = 3    /* sensor offline / out of range */
} EventState;
\end{lstlisting}

Design rules applied:
\begin{itemize}
  \item \textbf{Invalid readings never alert} (NF-07). A disconnected HC-SR04 reports \texttt{INVALID} but does not stop the temperature and motion checks.
  \item Only \texttt{STATE\ttu CRITICAL} triggers an \texttt{[ALERT]}. \texttt{STATE\ttu WARNING} is classified, but the current firmware prints no separate warning line.
  \item The motion metric is a \textbf{change from baseline}, not absolute acceleration, so it is orientation-independent.
\end{itemize}

%% ---------------------------------------------------------------- 9
\section{UART Alert Protocol}
Plain ASCII, one message per line, terminated by \texttt{\textbackslash r\textbackslash n}. Link parameters: \textbf{115200 baud, 8 data bits, no parity, 1 stop bit}.

\begin{lstlisting}
[INFO]  SYSTEM   | BOOT     | Vehicle Event Monitor v1.0
[INFO]  SYSTEM   | INIT     | Calibrating sensors...
[INFO]  SYSTEM   | INIT     | MPU6050 baseline calibrated
[WARN]  SYSTEM   | INIT     | MPU6050 not detected
[ALERT] MOTION   | CRITICAL | Ax=+1.82g Ay=-0.12g Az=+0.03g
[ALERT] EXTI     | CRITICAL | Hardware Emergency Interrupt fired on PB12!
[ALERT] TEMP     | CRITICAL | T=73C (threshold=60C)
[ALERT] OBSTACLE | CRITICAL | d=14cm (threshold=20cm)
[INFO]  SYSTEM   | NORMAL   | All sensors OK
[DATA]  T=26C | d=31cm | Ax=+0.01g Ay=-0.02g Az=+0.00g
\end{lstlisting}

\begin{tabularx}{\textwidth}{@{}L{3.8cm} X@{}}
\toprule
\textbf{Message} & \textbf{When}\\
\midrule
\texttt{BOOT} & Once at power-up (confirms UART is alive)\\
\texttt{INIT} & Startup; reports MPU6050 calibrated or not detected\\
\texttt{ALERT MOTION} & Impact of at least 1.5~g from baseline\\
\texttt{ALERT EXTI} & SOS button on PB12 pressed\\
\texttt{ALERT TEMP} & Over-temperature\\
\texttt{ALERT OBSTACLE} & Obstacle within 20~cm\\
\texttt{NORMAL} & Heartbeat about every 1~s while nothing is critical\\
\texttt{DATA} & Telemetry block for the dashboard (\texttt{d=--cm} if invalid)\\
\bottomrule
\end{tabularx}

%% ---------------------------------------------------------------- 10
\section{Testing and Validation}
\subsection{Standalone Test Programs}
Each test is a separate \texttt{main()} that can be flashed in place of \texttt{main.c}.
\begin{itemize}
  \item \textbf{\texttt{hrc\ttu test.c}} -- HC-SR04 bring-up. Counts valid versus total pings, exposes a stage diagnostic (\texttt{echo\ttu diag\ttu stage}: 1~trig sent, 2~echo high, 3~echo low, 4~OK, $-1$/$-2$~timeouts), toggles PC13 on each valid reading, and \textbf{auto-scans backup TRIG/ECHO pin pairs} (PB0/PB1, PA3/PA4, PA1/PA2 and swapped polarities) when no echo is seen.
  \item \textbf{\texttt{mpu6050\ttu test.c}} -- MPU6050 bring-up. \textbf{Auto-probes SCL/SDA pin pairs and I\textsuperscript{2}C addresses}, reads \texttt{WHO\ttu AM\ttu I}, computes the acceleration magnitude ($\approx$1.00~g at rest), and shows bus health on the LED (10~Hz when streaming, 1~Hz when searching).
\end{itemize}

\subsection{Test Plan (from the PRD)}
\begin{tabularx}{\textwidth}{@{}L{3.4cm} L{5.2cm} X@{}}
\toprule
\textbf{Test} & \textbf{Procedure} & \textbf{Pass criterion}\\
\midrule
LM35 normal & Read at room temperature & $\approx$25~\degC, \texttt{STATE\ttu NORMAL}\\
LM35 response & Warm the sensor by hand & Reading rises\\
LM35 critical & Temporarily lower \texttt{TEMP\ttu THRESHOLD\ttu C} (e.g.\ 30) & \texttt{[ALERT] TEMP} fires\\
MPU6050 identity & Read \texttt{WHO\ttu AM\ttu I} & Valid ID (PRD expects 0x68)\\
MPU6050 stationary & Leave still for 30~s & Deltas $\approx$0; no false trigger\\
MPU6050 impact & Tap or shake the board & \texttt{[ALERT] MOTION} within one reading\\
HC-SR04 accuracy & Object at a known distance (e.g.\ 30~cm) & Within $\pm$3~cm\\
HC-SR04 critical & Object at 10~cm & \texttt{[ALERT] OBSTACLE} fires\\
HC-SR04 timeout & Block or unplug ECHO & \texttt{STATE\ttu INVALID}, no alert\\
EXTI & Press PB12 button & \texttt{[ALERT] EXTI} printed, LED strobes\\
Integration & All sensors for 60~s & No spurious alerts; sensors alert independently\\
Power cycle & Reset the board & Banner printed, baseline re-established, no boot alerts\\
\bottomrule
\end{tabularx}

\subsection{Recorded Results}
\textit{To be completed by the team.} The repository contains no captured measurements, so no pass/fail results are claimed here. Record the bench results below before submission.

\begin{tabularx}{\textwidth}{@{}L{5cm} L{3.6cm} X L{1.8cm}@{}}
\toprule
\textbf{Test} & \textbf{Expected} & \textbf{Measured} & \textbf{Pass/Fail}\\
\midrule
LM35 at room temperature & $\approx$25~\degC & & \\
HC-SR04 at 30~cm & $30\pm 3$~cm & & \\
MPU6050 \texttt{WHO\ttu AM\ttu I} & 0x68 & & \\
MPU6050 stationary 30~s & No false alerts & & \\
Impact / shake & \texttt{[ALERT] MOTION} & & \\
SOS button (PB12) & \texttt{[ALERT] EXTI} & & \\
60~s integration run & No spurious alerts & & \\
\bottomrule
\end{tabularx}

%% ---------------------------------------------------------------- 11
\section{Implementation vs.\ Planned Design}
The firmware departs from the PRD in several places; they are listed so the report matches the code.

\begin{tabularx}{\textwidth}{@{}L{2.8cm} L{4.6cm} X@{}}
\toprule
\textbf{Topic} & \textbf{PRD / plan} & \textbf{Implemented}\\
\midrule
MCU part & Plan: STM32F401CCUx; PRD: STM32F411CEU6 & \textbf{STM32F411CEU6} (as in \texttt{config.h})\\
System clock & 84~MHz & \textbf{16~MHz HSI} (no PLL)\\
I\textsuperscript{2}C & Hardware I2C1 peripheral & \textbf{Bit-banged} I\textsuperscript{2}C on PB6/PB7 (more robust bring-up)\\
HC-SR04 capture & Timer input capture, TIM1\_CH2 & \textbf{GPIO polling with DWT cycle counter} on PB0/PB1\\
UART & USART2 on PA2 & USART1 (PA9), USART2 (PA2) and \textbf{two software UART pins}, all mirrored\\
Arithmetic & Integer only (NF-04) & Integer for the impact metric; \textbf{\texttt{float}} for temperature, distance and display values\\
Module layout & Separate \texttt{.c/.h} per driver & \textbf{Single \texttt{main.c}} plus \texttt{config.h}\\
Threshold levels & Single critical threshold & Added \textbf{Warning} tier and an \texttt{INVALID} state\\
Extra feature & --- & PB12 \textbf{EXTI emergency input}\\
Buzzer, RTC, SD & Buzzer in PRD (optional); RTC/SD in plan & Not implemented (Section~\ref{sec:outofscope})\\
\bottomrule
\end{tabularx}

%% ---------------------------------------------------------------- 12
\section{Limitations and Known Issues}
\begin{itemize}
  \item \textbf{Blocking, polled design.} Sensor reads and software UART transmission run in the super-loop. A bit-banged character takes about 87~\us, so a full alert line blocks the loop for several milliseconds, more than the 1~ms bound in PRD NF-02. The 100~ms period is still met.
  \item \textbf{Floating point in firmware} contradicts the PRD's integer-only rule (NF-04). It is harmless on the Cortex-M4F but would matter on an FPU-less port.
  \item \textbf{\texttt{snprintf} with \texttt{\%f}} increases code size and needs a suitable C library configuration in Keil.
  \item \textbf{Shared-variable access.} \texttt{state\ttu motion} and \texttt{critical\ttu alert\ttu count} are written from both the ISR and the main loop; the read-modify-write on the counter is not atomic.
  \item \textbf{PB12 handling.} The ISR forces \texttt{STATE\ttu CRITICAL}, but the next MPU6050 read overwrites \texttt{state\ttu motion}, so the critical state can be cleared before the alert is printed. The \texttt{emergency\ttu exti\ttu triggered} flag is only consumed while the state is still critical. There is also no debounce on the button.
  \item \textbf{MPU6050 DLPF of $\sim$5~Hz} smooths very short impacts, so a very brief shock could be attenuated before it reaches the threshold.
  \item \textbf{Alert rate limiting} (every 5th loop) means a critical state is first reported up to $\sim$0.5~s late.
  \item \textbf{No persistence or timestamps.} Events are not stored, so history is lost when the monitor is closed.
  \item \textbf{Documentation mismatch.} Header-pin comments in \texttt{mpu6050\ttu test.c} and \texttt{main.c} give different physical positions for PB6/PB7; verify against the board silkscreen.
  \item \textbf{No measured results} are recorded in the repository (Section~10.3).
\end{itemize}

%% ---------------------------------------------------------------- 13
\section{Future Work}
Implementation improvements to the existing system:
\begin{enumerate}
  \item Move to the \textbf{84~MHz PLL} clock and hardware \textbf{I\textsuperscript{2}C1 / timer input-capture} peripherals, as in the PRD, to free the CPU.
  \item Make UART transmission \textbf{interrupt- or DMA-driven} with a ring buffer so sensor reads are never delayed.
  \item Replace floating-point maths with integer scaling, and split \texttt{main.c} into per-driver \texttt{.c/.h} modules with unit tests.
  \item Make the PB12 path race-free: latch the event in a dedicated flag and debounce it in hardware or software.
  \item Add a collision-versus-rough-road classifier based on event duration and multi-axis signature, and a watchdog timer for fault recovery. These are planned enhancements, not completed features.
\end{enumerate}
Planned hardware and feature additions that were not built are listed in Section~\ref{sec:outofscope}.

%% ---------------------------------------------------------------- 14
\section{Repository Structure}
\begin{lstlisting}
Vehicle_Event_Logging/
|-- main.c          Application firmware (drivers + main loop)
|-- config.h        Pins, thresholds, EventState enum
|-- hrc_test.c      HC-SR04 standalone test
|-- mpu6050_test.c  MPU6050 standalone test
|-- PRD_Vehicle_Event_Logging.md
|-- Vehicle_Event_Logging_Response_System_Plan.md / .xlsx
|-- team11.pptx     Presentation slides
|-- REPORT.md / REPORT.tex / REPORT.pdf
\end{lstlisting}

%% ---------------------------------------------------------------- 15
\section{Build, Flash and Run}
\textbf{Toolchain:} Keil \textmu Vision with the STM32F4 device pack (device \texttt{STM32F411CEUx}), CMSIS device header \texttt{stm32f411xe.h}.
\begin{enumerate}
  \item Create a project for the STM32F411CEU6 and add \texttt{main.c} and \texttt{config.h}. To run a standalone test instead, add \texttt{hrc\ttu test.c} or \texttt{mpu6050\ttu test.c} in place of \texttt{main.c}.
  \item Enable the standard C library (\texttt{stdio}) so \texttt{snprintf} works.
  \item Build, then flash through ST-Link (SWD).
  \item Connect a USB-to-UART adapter RX to \textbf{PA9} (or PA2 / PB10 / PB4) and GND to GND.
  \item Open a terminal (Docklight / PuTTY) at \textbf{115200, 8-N-1}.
  \item Reset the board. Expect the boot banner, then the heartbeat and telemetry lines once per second.
\end{enumerate}

%% ---------------------------------------------------------------- 16
\section{Conclusion}
The project delivers a working three-sensor vehicle monitor on a bare-metal STM32F411: analog sensing (LM35 on the ADC), digital bus sensing (MPU6050 over a bit-banged I\textsuperscript{2}C master with bus recovery), timing-based sensing (HC-SR04 using the DWT cycle counter), an interrupt-driven emergency input, a four-way mirrored UART alert channel and local LED indication. It follows the PRD's classification model and message format and handles faults gracefully: invalid readings never raise alerts, and an offline MPU6050 is re-initialised automatically. The main gaps are the planned logging features, integer-only arithmetic and non-blocking communication. These, together with recorded bench measurements, are the clear next steps.

%% ---------------------------------------------------------------- 17
\section{Out of Scope (Planned but Not Implemented)}
\label{sec:outofscope}
The following features were considered in the original project plan but have not been implemented in the current firmware.

\subsection*{A. Data Logging and Timestamping}
\begin{itemize}
  \item DS3231 RTC module for real-time clock functionality.
  \item Timestamping of detected vehicle events.
  \item Micro-SD card integration using SPI.
  \item Persistent storage of event history.
\end{itemize}

\subsection*{B. Additional Sensor Integration}
\begin{itemize}
  \item INA219 voltage and current monitoring.
  \item Hall-effect sensor for vehicle speed measurement.
  \item LDR (Light Dependent Resistor).
  \item IR sensor.
  \item MAX30100 sensor.
\end{itemize}

\subsection*{C. User Interface and Indication}
\begin{itemize}
  \item OLED display for real-time system status.
  \item Buzzer for audible emergency alerts.
\end{itemize}

\subsection*{D. Wireless and Vehicle Communication}
\begin{itemize}
  \item Bluetooth communication.
  \item CAN (Controller Area Network) communication.
\end{itemize}

\subsection*{E. Vehicle Control and Safety Mechanisms}
\begin{itemize}
  \item Automatic braking mechanism.
  \item Motor actuation and control.
\end{itemize}

\subsection*{F. Advanced Processing}
\begin{itemize}
  \item Collision-versus-rough-road classification.
  \item Advanced event analysis using event duration and multi-axis signatures.
  \item Watchdog-based fault recovery.
\end{itemize}

\textit{Note:} The advanced-processing and watchdog items also appear in Section~13 (Future Work). They are planned enhancements, not completed features.

\vfill
\begin{center}\small\itshape
Prepared for 23CSE304 -- Embedded System, Department of Computer Science and Engineering,\\
School of Computing, Amrita Vishwa Vidyapeetham, Coimbatore, India. Mentor: Dr.\ A.\ Baskar.
\end{center}

\end{document}
